\chapter{研究对象与方法}

\section{研究对象}

研究对象由系统边界、输入、处理过程、输出和外部约束共同定义。验证开始前，应将自然语言目标转化为可检查的指标，并记录数据来源、版本和适用范围。

\subsection{验证原则}

验证遵循可复核性、可追踪性和一致性原则。关键结论应至少对应一项可定位证据，量和单位采用国家法定计量单位，术语、符号和缩略语在全文中保持一致。

\subsubsection{评价量}

设第 $i$ 次独立试验的观测值为 $x_i$，样本均值为
\begin{equation}
  \bar{x}=\frac{1}{n}\sum_{i=1}^{n}x_i .
  \label{eq:mean}
\end{equation}
式中，$n$ 为独立试验次数；$x_i$ 与 $\bar{x}$ 的计量单位相同。需要相互参照的公式才编号，编号置于行末圆括号内。

\section{实施流程}

验证流程包括需求分解、证据设计、试验执行、结果复核和报告归档五个阶段。每一阶段都应保留输入、输出、责任人和变更原因。

\begin{figure}[htbp]
  \centering
  \fbox{\begin{minipage}[c][32mm][c]{0.78\textwidth}
    \centering
    图形占位区：请替换为清晰的 PDF、PNG 或 EPS 文件。\\
    推荐优先使用可缩放的矢量图。
  \end{minipage}}\par
  \caption{验证流程示意}
  \label{fig:workflow}
\end{figure}

图\ref{fig:workflow}应紧随首次引用它的文字。图应具有自明性，图号和图题置于图的下方。
